\documentclass[10pt,a4paper]{amsart}
\usepackage[margin=19mm]{geometry}
\usepackage{amsmath,amssymb,mathtools}
\usepackage[scr=rsfs,cal=euler]{mathalfa}
\usepackage{xfrac}
\usepackage[unicode,hidelinks,bookmarks=false]{hyperref}
\newcommand{\ie}{i.e.\ }
\newcommand{\cf}{cf.\ }
\newcommand{\resp}{respectively\ }
\newcommand{\Subsection}{\subsection}
\providecommand{\nobreakdash}{\nobreak\hbox{-}}
\setlength{\emergencystretch}{2em}
\multlinegap0pt
\newtheorem{thm}{Theorem}
\usepackage{cleveref}
\crefname{thm}{Theorem}{Theorems}
\newcommand{\Brac}[1]{\left(#1\right)}
\newcommand{\Gap}{\mathsf{Gap}}
\newcommand{\Rectbrac}[1]{\left[#1\right]}
\newcommand{\St}{\Rightarrow}
\newcommand{\TV}{\mathsf{TV}}
\newcommand{\mR}{\mathbb{R}}
\newcommand{\matc}{^{\rm c}}
\newcommand{\norme}[1]{\left|#1\right|}
\newcommand{\sfP}{\mathsf{P}}
\title{Spectral Gap of Metropolis Algorithms for Non-smooth Distributions under Isoperimetry}
\author{Shuigen Liu, Xin T. Tong}
\date{Source: arXiv:2601.21763v2 (CC BY 4.0)}
\begin{document}
\maketitle

\noindent\emph{Source and licence.} Spectral Gap of Metropolis Algorithms for Non-smooth Distributions under Isoperimetry. Version arXiv:2601.21763v2 (CC BY 4.0), distributed by arXiv under the Creative Commons Attribution 4.0 International licence, http://creativecommons.org/licenses/by/4.0/, as recorded in the arXiv licence metadata for this exact version and shown on the abstract page https://arxiv.org/abs/2601.21763v2. Full licence notice: \texttt{licenses/arxiv-2601.21763-CC-BY-4.0.txt}.\par\smallskip

\noindent\textit{Editorial scope.} Counted unit: the article's thm thm:ChIneq (source0.tex lines 335--343). The article's complete original proof of it is delivered with it (source0.tex lines 951--1018, one \texttt{proof} environment of 68 lines, which the article places away from the statement). No mathematics is written, repaired or re-derived: the statement and the proof are the article's own lines, in the article's own notation, and no retained line is substituted or re-flowed.  Retained uncounted as the source-local context the proof needs: lines 329--332; lines 374--377; lines 416--419.  Nothing was omitted from any retained span, so the package carries no omission record.  The retained lines cite 1 works, whose entries are transcribed verbatim from the article's own bibliography source in the e-print and delivered with short editorial labels recorded in the item's reference mapping.  No retained line contains a figure environment, an image-inclusion command or drawing code, so no image is delivered and the package declares no asset dependency.  Editorial formatting only, in addition: an AMS \texttt{amsart} preamble that restates the article's own declarations and macros used by the retained lines, namely the environments \texttt{thm} on separate counters, together with the standard packages those lines need.  The retained environments print in this document as Theorem 1 in order of appearance, whereas the article prints the same items with the numbering of its own full document.  The complete native source of the article as distributed in the arXiv e-print for this exact version is bundled unchanged as \texttt{sources/193611/source0.tex}.
\begin{equation}
\label{eqn:Cond}
    \Phi(\sfP) = \inf \left\{ \, \frac{ \Brac{\pi \otimes \sfP } (S \times S\matc)}{\pi(S)} \, :\; S \subseteq \mR^d,~ 0 < \pi(S) \le \frac{1}{2} \right\}.
\end{equation}

    \begin{equation}
    \label{eqn:CloseCouple}
        \norme{ x-y } < \delta ~\St~ \TV (\delta_x \sfP,\delta_y \sfP) < 1-\epsilon .
    \end{equation}

\begin{equation}
\label{eqn:SpecGap_via_Iso}
    \Gap(\sfP) \ge \frac{\epsilon^2}{8} \Rectbrac{ \sup_{\theta\in[0,1]} \min \Big\{ 1-\theta , \Upsilon(\delta) \inf_{t\in(0,\frac{1}{2}]} \frac{ F(\theta t) }{ 2t } \Big\} }^2.
\end{equation}

\begin{thm}[\cite{MR930082}]
\label{thm:ChIneq}
Let $\sfP$ be a reversible Markov kernel. Then it holds
\begin{equation}
\label{eqn:ChIneq}
    \frac{\kappa}{2} \Phi(\sfP)^2 \le \Gap(\sfP) \le \Phi(\sfP),
\end{equation}
where $\kappa \ge 1$ is a universal constant.
\end{thm}

\begin{proof}
By Cheeger's inequality \Cref{thm:ChIneq}, it suffices to lower bound the Cheeger constant $\Phi(\sfP)$ defined in \eqref{eqn:Cond}.
%
For any measurable set $S \subseteq \mR^d$, denote the three sets
\[
    S_1 = \left\{ x \in S: \sfP(x,S\matc) \le \epsilon/2 \right\}, \quad S_2 = \left\{ x \in S\matc: \sfP(x,S) \le \epsilon/2 \right\}, \quad S_3 = (S_1\cup S_2)\matc.
\]
By definition, for any $x\in S_1$ and $y \in S_2$,
\[
    \TV \Brac{ \delta_x \sfP , \delta_y \sfP } \ge \sfP(x,S) - \sfP(y,S) \ge 1 - \epsilon/2 - \epsilon/2 = 1 - \epsilon.
\]
By close coupling condition \eqref{eqn:CloseCouple}, this implies that $\norme{x-y} \ge \delta$, and thus $d(S_1,S_2) \ge \delta$. 


For any $\theta \in (0,1)$, we discuss by the following two cases: 
\begin{itemize}
    \item Case I: $\pi(S_1) \le \theta \pi(S)$ or $\pi(S_2) \le \theta \pi(S\matc)$.
    \item Case II: $ \pi(S_1) > \theta \pi(S)$ and $ \pi(S_2) > \theta \pi(S\matc)$.
\end{itemize}
(1) Denote for simplicity $\Pi = \pi \otimes \sfP$. If $\pi(S_1) \le \theta \pi(S)$, by definition of $S_1$, we have
\[
    \Pi(S \times S\matc) \ge \Pi((S\setminus S_1) \times S\matc) \ge \frac{\epsilon}{2}  \pi(S\setminus S_1) \ge \frac{\epsilon}{2} (1-\theta) \pi(S).
\]
%
Since $\sfP$ is reversible with respect to $\pi$, we have $\Pi(S \times S\matc) = \Pi(S\matc \times S)$, and thus the above inequality also holds for $\Pi(S\matc \times S)$.
Likewise, if $\pi(S_2) \le \theta \pi(S\matc)$, we can show that
\[
    \Pi(S\matc \times S) \ge \frac{\epsilon}{2} (1-\theta) \pi(S\matc).
\]
Therefore, in Case I, it always holds that
\[
    \Pi(S \times S\matc) \ge \frac{\epsilon}{2} (1-\theta) \min \{ \pi(S), \pi(S\matc) \}.
\]
(2) When $ \pi(S_1) > \theta \pi(S)$ and $ \pi(S_2) > \theta \pi(S\matc)$, by the three-set isoperimetric inequality,
\[
    \pi(S_3) \ge \Upsilon \Brac{ d(S_1,S_2) } F \Brac{ \min \{ \pi(S_1), \pi(S_2) \} } \ge \Upsilon(\delta)\, F \Brac{ \theta \min \{ \pi(S), \pi(S\matc) \} }.
\]
where the second inequality follows from $d(S_1,S_2) \ge \delta$ and the monotonicity of $\Upsilon$ and $F$.
Now by the definition of $S_1,S_2$, we have
\begin{align*}
    \Pi(S \times S\matc) =~& \frac{1}{2} \Brac{ \Pi(S \times S\matc) + \Pi(S\matc \times S) } \\
    \ge~& \frac{1}{2} \Brac{ \Pi((S\setminus S_1) \times S\matc) + \Pi((S\matc\setminus S_2) \times S) } \\
    \ge~& \frac{\epsilon}{4} \Brac{ \pi(S\setminus S_1) + \pi(S\matc\setminus S_2)  } = \frac{\epsilon}{4} \pi(S_3) \\
    \ge~& \frac{\epsilon}{4} \Upsilon(\delta) F \Brac{ \theta \min \{ \pi(S), \pi(S\matc) \} }.
\end{align*}
%
Combining the above two cases, we have
\[
    \Pi(S \times S\matc) \ge \min \left\{ \frac{\epsilon}{2} (1-\theta) \min \{ \pi(S), \pi(S\matc) \}, \frac{\epsilon}{4} \Upsilon(\delta) F \Brac{ \theta \min \{ \pi(S), \pi(S\matc) \} } \right\}.
\]
By the symmetry of $S$ and $S\matc$, we can assume without loss of generality that $\pi(S) \le \frac{1}{2}$.
Then
\[
    \frac{\Pi(S \times S\matc)}{\pi(S)} \ge \min \left\{ \frac{\epsilon}{2} (1-\theta), \frac{\epsilon}{4} \Upsilon(\delta) \frac{ F(\theta \pi(S)) }{ \pi(S) } \right\}.
\]
Taking infimum over $S$, we get
\[
    \Phi(\sfP) \ge \frac{\epsilon}{2} \min \left\{ 1-\theta , \Upsilon(\delta) \inf_{t\in(0,\frac{1}{2}]} \frac{ F(\theta t) }{ 2t } \right\}.
\]
The first claim follows by taking supremum over $\theta\in[0,1]$ and applying \eqref{eqn:ChIneq}.

For the second claim, when $F$ is a concave function on $(0,\frac{1}{2}]$, then $\forall 0<x\le y\le \frac{1}{2}$,
\[
    \frac{F(y)-F(x)}{y-x} \le \frac{F(x) - \lim_{t \to 0+} F(t) }{x - 0} \le \frac{F(x)}{x} ~\St~ \frac{F(y)}{y} \le \frac{F(x)}{x}.
\]
Therefore, $\frac{F(t)}{t}$ is non-increasing on $(0,\frac{1}{2}]$, and thus $\inf_{t\in(0,\frac{1}{2}]} \frac{ F(\theta t) }{ 2t } = F(\theta/2)$.
Taking $\theta = \frac{1}{2}$ in \eqref{eqn:SpecGap_via_Iso} yields the desired result.
\end{proof}

\begin{thebibliography}{99}
\bibitem[MR9300]{MR930082} Lawler, G.~F. and Sokal, A.~D. (1988). \newblock Bounds on the {$L^2$} spectrum for {M}arkov chains and {M}arkov processes: a generalization of {C}heeger's inequality. \newblock {\em Trans. Amer. Math. Soc.}, 309(2):557--580.
\end{thebibliography}
\end{document}
